\section{La Meta (Capítulos 16–26): V/F Posibles                                                           48}

\subsection{Conceptos clave Cap. 16–26 . . . . . . . . . . . . . . . . . . . . . . . . . . . . . .      48}

\subsection{20+ Afirmaciones V/F sobre La Meta (Cap. 16–26) . . . . . . . . . . . . . . . . .           48}
\section{La Meta (Capítulos 16–26): V/F Posibles}

\begin{tipprueba}{Instrucciones para la parte digital}

   Si la afirmación es Falsa, debes argumentar brevemente por qué (1–2 oraciones). Los
   capítulos 16–26 cubren el Proceso de Mejora Continua (POOGI), los cinco pasos
   de focalización, el tambor-amortiguador-cuerda (DBR), el conflicto con Hilton
   Smith y el cierre de la trama (Alex es ascendido).


\end{tipprueba}

\subsection{Conceptos clave Cap. 16–26}

\begin{definicion}{Los cinco pasos de focalización (Cap. 16–18)}

   El proceso continuo de mejora de la Teoría de Restricciones (TOC):

1. Identificar la restricción (cuello de botella).
2. Explotar la restricción (sacarle el máximo sin invertir).
\end{definicion}
\section{Subordinar todo lo demás a la restricción.}
4. Elevar la restricción (invertir capacidad).
5. Volver al paso 1 (evitar la inercia): rota una restricción, nace otra.

\begin{definicion}{Tambor-Amortiguador-Cuerda (DBR) (Cap. 17–19)}

   Tambor (Drum): el cuello de botella marca el ritmo de toda la planta.
   Amortiguador (Buffer): inventario de protección antes del CB para que nunca se
   quede sin trabajo (protege el throughput contra la variabilidad).
   Cuerda (Rope): señal que ata la liberación de material al ritmo del tambor, evitando
   exceso de WIP.

\end{definicion}

\begin{definicion}{El conflicto y el cierre (Cap. 20–26)}

   Hilton Smith (rival burocrático) presiona con métricas de eficiencia local; Alex de-
   fiende el throughput global.
   Alex y su equipo (Lou, Bob, Stacey, Ralph) generalizan el método y descubren que
   la gestión es un proceso de preguntar: ¿Qué cambiar? ¿Hacia qué cambiar? ¿Cómo
   provocar el cambio? (Procesos de Pensamiento).
   El éxito de la planta lleva a Alex a ser ascendido a director de división.
   Lección final: las habilidades del científico (método socrático de Jonah) son la
   verdadera herramienta de gestión.


\end{definicion}

\subsection{20+ Afirmaciones V/F sobre La Meta (Cap. 16–26)}
#RAfirmación y justificación
1VEl primer paso de TOC es identificar la restricción del sistema.
  Sin saber dónde está el CB, mejorar otra cosa es desperdicio.
#RAfirmación y justificación
2F“Explotar” la restricción significa invertir dinero para ampliarla.
  Falso: explotar es sacarle el máximo SIN invertir; invertir es “elevar”
  (paso 4).
3VEl amortiguador (buffer) se coloca antes del cuello de botella.
  Protege al CB de quedarse sin material por la variabilidad aguas arriba.
4FLa “cuerda” libera material al máximo de la capacidad de la primera
  estación.
  Falso: la cuerda ata la liberación al ritmo del TAMBOR (el CB), no de
  la primera estación.
5VSubordinar significa que los recursos no-cuello de botella trabajan al
  ritmo del CB.
  Todo se sincroniza con el tambor; el ocio en no-CB es sano.
6FUna vez elevada la restricción, el proceso de mejora termina.
  Falso: hay que volver al paso 1 (evitar la inercia); aparece una nueva
  restricción.
7VAl romper una restricción, ésta puede trasladarse a otra parte del sistema
  (incluso al mercado).
  Por eso el paso 5 obliga a reiniciar el ciclo.
8FHilton Smith representa la visión de throughput global que defiende Alex.
  Falso: Hilton encarna las métricas de eficiencia LOCAL que Alex com-
  bate.
9VEl tambor (drum) es el cuello de botella que fija el ritmo de toda la
  planta.
  La planta no puede producir más rápido que su restricción.
10   VReducir el tamaño de lote en la planta mejora el flujo y reduce el tiempo
  de ciclo.
  Lotes pequeños bajan el WIP (Ley de Little), aunque aumenten los se-
  tups.
11   FLa gestión efectiva consiste en dar respuestas y órdenes directas al equi-
  po.
  Falso: Jonah enseña el método socrático — preguntar para que el equipo
  descubra.
12   VLos “Procesos de Pensamiento” responden: qué cambiar, hacia qué cam-
  biar y cómo.
  Es la generalización del método al final del libro.
13   FMantener ocupados todos los recursos al 100 % maximiza el throughput.
  Falso: activar un no-CB sin que el CB procese solo acumula inventario
  (WIP).
14   VUn minuto perdido en el cuello de botella es un minuto perdido para
  todo el sistema.
  El CB determina el throughput; su tiempo es irrecuperable.
15   FUn minuto ahorrado en un recurso no-cuello de botella es un minuto
  ganado para el sistema.
  Falso: es un espejismo; solo genera más inventario antes del CB.
16   VEl éxito de la planta de Bearington lleva a Alex a ser ascendido a la
  división.
  Cierre de la trama: el método TOC se valida y se expande.
17   FEl buffer del CB debe ser lo más grande posible para máxima seguridad.
  Falso: un buffer excesivo eleva el WIP y el tiempo de ciclo; debe ser el
  justo.
18   VLa meta de la empresa sigue siendo ganar dinero ahora y en el futuro.
  Se mantiene de la primera mitad: throughput, inventario y gasto opera-
  tivo.
19   FSubordinar implica que el cuello de botella se adapta al ritmo de los
  demás recursos.
  Falso: es al revés — los demás se subordinan al ritmo del CB.
20   VLa inercia es el mayor enemigo del proceso de mejora continua.
#RAfirmación y justificación
  Las reglas creadas para una restricción quedan obsoletas cuando ésta
  cambia.
21   VEl equipo descubre que las habilidades de gestión son las del método
  científico.
  Observar, hipotetizar, verificar — como hace Jonah.
22   FDBR busca maximizar la eficiencia individual de cada máquina de la
  línea.
  Falso: busca sincronizar el flujo con el tambor, no la eficiencia local.
